\section{Introduction}
\label{sec:introduction}

The Standard Model and general relativity are extraordinarily successful at organizing experimental facts.
At the same time, their combined fundamental ``catalog''---multiple fields, symmetry sectors, symmetry breaking
structure, and many parameters fixed by measurement---invites a complementary line of inquiry:
whether a smaller set of ground ingredients could underlie this effective complexity.
This paper explores such a \emph{minimal-ingredient reconstruction} approach.

\paragraph{Substrate-first strategy.}
We adopt a deterministic \emph{substrate-first} strategy: posit an ontically real wave medium and treat familiar
field theories as coarse-grained descriptions of constrained wave dynamics.
The goal is \emph{not} to dispute the established phenomenology of quantum field theory, but to ask what could make
gauge structure, relativistic kinematics, and particle-like quantization appear naturally as emergent regularities.

\paragraph{Ground ingredients (minimal commitments).}
The theory fragment developed here aims to keep its primitive commitments small:
(i) a single continuous substrate with local action dynamics,
(ii) geometric nonlinearity arising from the embedding (rather than adding separate interaction fields by hand), and
(iii) a controlled separation between a linear wave sector and a nonlinear self-guidance sector.
Symmetry properties such as large-scale isotropy are treated as \emph{constraints on admissible constitutive choices}
(e.g., writing the energy in terms of invariants of the induced metric), not as additional ``ingredients''.
Everything else is treated as emergent bookkeeping: in particular, gauge potentials are interpreted as connection
data of an ordered wave sector, not as independent fundamental substances.

\paragraph{Substrate hypothesis.}
We assume a three-dimensional elastic medium (``brane'') embedded in $\R^4$, described by an embedding map
\begin{equation}
  \vecX:\Omega\times\R \to \R^4, \qquad (x,t)\mapsto \vecX(x,t), \qquad \Omega\subset\R^3,
  \label{eq:embedding-map}
\end{equation}
with an external time parameter $t$.
The dynamics is specified by a local hyperelastic action on $\Omega$, constructed from invariants of the
induced metric of the embedding. This is the minimal coordinate-free way to write an elastic energy without
committing the effective continuum description to a preferred direction; microscopically, a cubic lattice
substrate is allowed, and large-scale isotropy is treated as an empirical constraint on admissible lattice
couplings and parameter regimes \citep{Ogden1984,Holzapfel2000,DoCarmo1992}.

\paragraph{Gravity from transverse excitation (contraction).}
A central motivating idea---already present at the level of the induced metric---is that excitation in the
embedding dimension generically feeds back into \emph{lateral} stresses. Intuitively, if a region of the brane
acquires a transverse profile $X^4=u(x,t)$, then intrinsic line elements become longer by a Pythagorean
amount; in a pre-tensioned medium this extra arc length is paid for by additional tensile stress that pulls
neighboring material points inward. Coarse-grained over scales large compared to the carrier wavelength, such
slowly varying contraction fields can bias the propagation of other wave packets and the equilibrium of
localized modes. We therefore treat Newtonian gravity as an emergent, long-range \emph{contraction channel}
mediated by transverse excitation, not as a separately postulated fundamental field.

\paragraph{Narrowband order and geometric phase.}
A further postulate is the existence of an ordered, narrowband carrier sector that provides
a coherent phase reference over macroscopic scales (assumed to live primarily in a longitudinal
polarization branch).
Geometric phase and holonomy are universal in wave systems and in adiabatic transport of polarized states
\citep{Pancharatnam1956,Berry1984,WilczekZee1984,Bliokh2015SpinOrbit,Cisowski2022}.
Here they are promoted to a structural bridge: adiabatic transport of a polarization subspace induces a
Berry (rank-1) or Wilczek--Zee (non-Abelian) connection, which is interpreted as the natural origin of
gauge degrees of freedom in an effective description.

\paragraph{Particles as solitons; probability not fundamental.}
We do not take the probabilistic interpretation of quantum theory as fundamental.
Instead, apparent point-like interactions and discrete outcomes are attributed to deterministic wave dynamics
combined with Fourier-localization constraints and threshold-like nonlinear exchange events.
Particle-like states are localized solitons whose spatial structure is organized by spherical symmetry and
spherical-harmonic mode families.
The spin-$\tfrac12$ $4\pi$ aspect is attributed to geometric phase/holonomy of an internal polarization
subspace (a ``spinorial'' parallel transport effect) rather than to literal multi-winding of a photon path.

\paragraph{Relation to existing ideas.}
The emergence of an effective light cone and Lorentz-like kinematics from a preferred-time medium has
close analogues in analogue-gravity systems \citep{Barcelo2011AnalogueGravity,Volovik2003HeliumDroplet,Jacobson2004EinsteinAether}.
The present use of the term ``brane'' is literal (an elastic medium), and should not be confused with
field-theoretic brane-world scenarios in high-energy theory \citep{RandallSundrum1999}.

\paragraph{Paper organization.}
\Cref{sec:scope} states scope, contributions, and non-claims.
\Cref{sec:continuum-model} formulates the continuum substrate model (written in a coordinate-free form with an isotropic constitutive baseline).
\Cref{sec:emergent-relativity} discusses the linear-wave regime and an emergent effective light cone.
\Cref{sec:geophase,sec:effective-field} develop the geometric connection picture (Berry/Wilczek--Zee) and its
gauge-covariant envelope formulation.
\Cref{sec:localized} treats localized modes and shows how to pose them as self-adjoint eigenproblems around
symmetric backgrounds.
Discrete implementation details are collected in \Cref{app:discrete}.
